% =====================================================================
% 01 — Introduction.
% =====================================================================
% Planning (% comments only):
% Six-move expanded:
%   1. Problem motivation (slogans on both sides; we want a map)
%   2. Set up assumption (the bundling)
%   3. Flip (three independent indictments; layer discipline)
%   4. Instantiate (which claims survive)
%   5. Evaluation (the companion probe suite + literature)
%   6. Implications (where to invest research effort)
% =====================================================================

\section{Introduction}
\label{sec:introduction}

The discourse around autoregressive language models has settled into two slogans.
One says they work, scale, and only need more compute, data, and tools.
The other says they are fundamentally broken because each generated token has a nonzero probability of being wrong, and a length-$T$ generation is fully correct with probability at most $(1-\varepsilon)^T$, which decays exponentially.

Both slogans are too clean.

The first ignores that local fluency is not the same as global correctness on tasks like proof writing, program synthesis, or long-horizon planning.
The second ignores that the math is valid only under three premises: independent token errors, a constant per-token error rate, and unrecoverability of any single wrong token.
Modern deployed systems are designed precisely to violate those premises.
Tools catch syntax, arithmetic, and type errors;
verifiers reject invalid proofs;
search lets a model retry;
and reinforcement learning from verifier feedback shifts the empirical error process away from the i.i.d. assumption the bound requires.

Most critiques of autoregressive language models bundle three distinct design choices into a single object.
The first is the \emph{autoregressive factorization} $p(x) = \prod_t p(x_t \mid x_{<t})$, which trades the global partition function of a sequence-level model for $T$ local normalizers.
The second is the \emph{softmax parameterization} of each conditional, which factors as an exponential over a learned hidden state followed by normalization over the vocabulary.
The third is the \emph{maximum-likelihood objective with teacher forcing}, which trains each conditional on clean prefixes drawn from data rather than from the model's own past samples.
These three choices are independent.
Different objections indict different choices, and any honest case for an alternative architecture must specify which of the three it is replacing.

The contribution of this paper is a layer-disciplined map of the standard objections to autoregressive language models.
We tag each objection by the layer at which it lives.
\emph{Architecture}-layer claims (\eg the softmax bottleneck, pure-attention rank collapse, fixed compute per token) can be tested on a random-initialized model.
\emph{Objective}-layer claims (\eg mode-covering pressure from forward-KL minimization) require a controlled training run.
\emph{Trained-behavior}-layer claims (\eg error compounding, reversal failures, hallucination rates) require both a trained model and an evaluator with a notion of correctness.
\emph{External}-layer claims (\eg the data wall) are properties of a resource curve, not of any single model.
Mixing layers is the most common way to fool oneself about which mechanism is doing the work.

We then ask, for each objection, whether it is a theorem, a robust empirical finding, a weak empirical finding, or refuted under modern conditions.
Two clusters emerge.
The softmax bottleneck~\citep{yang2018softmax}, doubly-exponential rank collapse in pure attention~\citep{dong2021attention}, and forward-KL mode-covering pressure~\citep{theis2016note,minka2005divergence} are theorems whose preconditions are visible from the parameterization or the objective.
The error-compounding bound~\citep{lecun2022autonomous,bachmann2024pitfalls}, the reversal curse~\citep{berglund2023reversal}, and most hallucination claims are empirical patterns whose realized effect under verifier-augmented systems is an open quantitative question rather than a fundamental impossibility.

Alternatives to the autoregressive recipe inherit, evade, or relocate the same costs.
Energy-based models~\citep{lecun2006tutorial,grathwohl2019your} score whole objects globally and avoid the local-then-global mismatch, but pay for the partition function or its approximations at training and at inference.
Score matching~\citep{hyvarinen2005score}, noise-contrastive estimation~\citep{gutmann2010noise}, and ratio matching~\citep{hyvarinen2007some} sidestep $Z_\theta$ during training but typically push the cost to inference-time sampling.
Diffusion and denoising models~\citep{sohldickstein2015deep,ho2020denoising,song2021scorebased,nichol2021improved} score a hierarchy of perturbations to the data and yield a tractable training loss at the cost of an iterative inference loop.
Discrete and masked diffusion~\citep{lou2023discrete,nie2024scaling} bring the same idea to vocabulary-indexed sequences.
Joint-embedding predictive architectures~\citep{lecun2022autonomous} replace per-token reconstruction with feature-space prediction.
Each is a different choice about \emph{where the hard part lives}.

This paper is the literature-grounded half of a two-part project.
The companion experiment, \texttt{experiments/02\_ar\_pros\_cons} in the same repository, tests each surviving claim with an isolated probe at the correct layer and then re-measures the same mechanism inside an end-to-end run on a domain with a hard verifier (Lean 4, via \textsc{VeriBench}).
The review-paper coordination layer, \texttt{experiments/03\_review\_paper}, keeps the paper, blog series, and experiment plan synchronized and adds a simple vision track, initially \textsc{MNIST}, before moving to a richer dataset.
We refer to that suite throughout, both for which claims it pre-registers as load-bearing and for which it pre-registers as likely masked once everything is trained together.

The remainder of the paper is organized as follows.
\Cref{sec:advantages} restates the honest case for autoregressive language models: the four design properties that explain why they scaled.
\Cref{sec:catalog} catalogs the standard objections, with layer tags.
\Cref{sec:what_holds} evaluates each objection against the bar of theorem, robust empirical, weak empirical, or refuted, and identifies the three independent indictments the surviving claims support.
\Cref{sec:experimental_program} states the companion experimental program: toy controls, \textsc{VeriBench}/Lean, and an initial \textsc{MNIST} vision track.
\Cref{sec:ebm_motivation} introduces energy-based models from the global-scoring viewpoint and isolates the partition-function obstacle as the central cost.
\Cref{sec:ebm_training,sec:ebm_inference} survey the training and inference methods that aim to evade or amortize that cost.
\Cref{sec:other_alts} surveys non-EBM alternatives.
\Cref{sec:conclusion} discusses where the hard part should live and the empirical questions that this paper deliberately defers to the companion suite.
